\documentclass[a4paper,12pt]{article}
\usepackage[utf8]{inputenc}
\usepackage[ngerman]{babel}
\usepackage{amsmath}
\usepackage{amssymb}
\usepackage{amsthm}
\usepackage{geometry}
\usepackage{graphicx}
\usepackage[hidelinks]{hyperref}
\usepackage{listings}
\usepackage{xcolor}
\usepackage{enumitem}
% bbm and mathrsfs not available in this environment

\geometry{margin=2.5cm}

\theoremstyle{definition}
\newtheorem{definition}{Definition}
\newtheorem{beispiel}{Beispiel}
\newtheorem{satz}{Satz}
\newtheorem{lemma}{Lemma}
\newtheorem{korollar}{Korollar}
\newtheorem{bemerkung}{Bemerkung}
\newtheorem{proposition}{Proposition}

\setlist[itemize,1]{label=--}

% ---------- Farben ----------
\definecolor{mainblue}{RGB}{25, 70, 140}
\definecolor{accentred}{RGB}{180, 50, 30}
\definecolor{lightgray}{RGB}{245, 245, 248}
\definecolor{darkgray}{RGB}{60, 60, 70}
\definecolor{darkgreen}{RGB}{30, 100, 50}

\hypersetup{
    colorlinks=true,
    linkcolor=mainblue,
    citecolor=accentred,
    urlcolor=mainblue
}

\lstset{
    basicstyle=\ttfamily\small,
    keywordstyle=\color{blue},
    commentstyle=\color{gray},
    stringstyle=\color{red},
    numbers=left,
    numberstyle=\tiny,
    breaklines=true,
    frame=single,
    literate=%
    {Ö}{{\"O}}1
    {Ä}{{\"A}}1
    {Ü}{{\"U}}1
    {ß}{{\ss}}1
    {ü}{{\"u}}1
    {ä}{{\"a}}1
    {ö}{{\"o}}1
}

\newcommand{\Z}{\mathbb{Z}}
\newcommand{\R}{\mathbb{R}}
\newcommand{\C}{\mathbb{C}}
\newcommand{\N}{\mathbb{N}}
\newcommand{\Lat}{\mathcal{L}}
\newcommand{\B}{\mathbf{B}}
\newcommand{\norm}[1]{\left\|#1\right\|}
\newcommand{\innerprod}[2]{\left\langle #1,\, #2 \right\rangle}

\title{\textbf{Gitterbasierte Kryptographie und Quanten-Fehlerkorrektur:}\\
\large Kürzeste-Vektor-Problem, GKP-Codes und Anonyme Reputationssysteme}
\author{Stephan Epp}
\date{3. Juni 2026}

\begin{document}

\maketitle

\begin{abstract}
Diese Arbeit gibt eine formale und ausführliche Einführung in drei eng verwandte Gebiete der
modernen Gittertheorie und quantensicheren Kryptographie: das Shortest-Vector-Problem~(SVP),
Gottesman-Kitaev-Preskill-Codes~(GKP) sowie gitterbasierte anonyme Reputationssysteme.
Sämtliche zentralen Ergebnisse werden rigoros bewiesen, durch matplotlib-Abbildungen
veranschaulicht und in den aktuellen Forschungsstand eingeordnet.
Der Subgraph Algorithmus~\cite{Epp2026Subgraph} liefert dabei die algorithmische Vergleichsgrundlage.
\end{abstract}

\tableofcontents
\newpage

% ============================================================
\section{Einführung}
% ============================================================

Die Sicherheit heutiger Public-Key-Kryptosysteme (RSA, ECDH, DSA) basiert auf
der vermuteten Schwierigkeit des diskreten Logarithmus und der Faktorisierung großer ganzer Zahlen.
Shors Algorithmus~\cite{Shor1994} zeigt, dass ein hinreichend großer Quantencomputer beide Probleme
in Polynomialzeit löst und damit diese Systeme vollständig bricht.

Gitterprobleme -- insbesondere das Shortest-Vector-Problem~(SVP) und das Closest-Vector-Problem~(CVP)
-- gelten bislang als resistent gegen Quantenangriffe und bilden die Basis einer post-quanten-sicheren
Kryptographie~\cite{Ajtai1996,Micciancio2009}.
Parallel dazu bieten Gottesman-Kitaev-Preskill-Codes (GKP)~\cite{GKP2001} einen
vielversprechenden Ansatz zur Quantenfehlerkorrektur, der direkt auf Gittern in der
symplektischen Phasenraumdarstellung basiert.

\subsection{Problemstellung und Motivation}

Diese Arbeit behandelt drei miteinander verknüpfte Fragestellungen:

\begin{enumerate}
    \item \textbf{SVP und Gitterbasisreduktion:}
          Wie schwer ist es, den kürzesten Vektor in einem Gitter $\Lat \subset \R^n$ zu finden?
          Welche Algorithmen existieren, und wie sind deren Laufzeiten?

    \item \textbf{GKP-Codes:}
          Wie lassen sich logische Qubits durch Gitterpunkte im Phasenraum kodieren?
          Welche Verbindung besteht zwischen Gittergeometrie und Codeabstand?

    \item \textbf{Anonyme Reputationssysteme:}
          Wie kann Anonymität in kryptographischen Review-Systemen durch gitterbasierte
          Zero-Knowledge-Beweise garantiert werden?
\end{enumerate}

\subsection{Beitrag dieser Arbeit}

\begin{itemize}
    \item Formale Beweise aller zentralen Aussagen zu SVP, GKP-Codedistanz und
          Systemsicherheit anonymer Reputationssysteme.
    \item Experimentelle Illustration aller Konzepte durch acht \texttt{matplotlib}-Abbildungen.
    \item Ausführlicher Ausblick mit über 20 offenen Forschungsrichtungen.
\end{itemize}

% ============================================================
\section{Gittertheoretische Grundlagen}
% ============================================================

\subsection{Gitter und Basen}

\begin{definition}[Gitter]
\label{def:gitter}
Seien $\mathbf{b}_1, \ldots, \mathbf{b}_n \in \R^m$ linear unabhängige Vektoren.
Das durch diese \emph{Basis} $\B = [\mathbf{b}_1 \mid \cdots \mid \mathbf{b}_n]$ erzeugte
\emph{Gitter} ist
\[
    \Lat(\B) = \left\{ \sum_{i=1}^{n} z_i \mathbf{b}_i \;\middle|\; z_i \in \Z \right\}
             = \{ \B \mathbf{z} \mid \mathbf{z} \in \Z^n \}.
\]
Man nennt $n$ den \emph{Rang} (oder die \emph{Dimension}) und $m \geq n$ die
\emph{Einbettungsdimension}.
Im Vollrang-Fall gilt $m = n$.
\end{definition}

\begin{bemerkung}
Die Basis ist \emph{nicht} eindeutig: Für jede unimodulare Matrix $\mathbf{U} \in \Z^{n \times n}$
mit $|\det \mathbf{U}| = 1$ erzeugen $\B$ und $\B \mathbf{U}$ dasselbe Gitter.
\end{bemerkung}

\begin{definition}[Fundamentalparallelotop]
Das \emph{Fundamentalparallelotop} einer Basis $\B$ ist
\[
    \mathcal{P}(\B) = \left\{ \B \mathbf{x} \;\middle|\; \mathbf{x} \in [0,1)^n \right\}.
\]
Es gilt $\det \Lat(\B) := |\det \B| = \mathrm{Vol}(\mathcal{P}(\B))$.
\end{definition}

\begin{definition}[Spektralnorm einer Basis]
Die \emph{Spektralnorm} $s_1(\B)$ ist der größte singuläre Wert von $\B$, also
\[
    s_1(\B) = \norm{\B}_2 = \sigma_{\max}(\B).
\]
Eine Basis mit kleiner Spektralnorm heißt \emph{gut} oder \emph{reduziert}.
\end{definition}

\begin{figure}[ht]
    \centering
    \includegraphics[width=0.75\textwidth]{plots/plot1_lattice_basis.pdf}
    \caption{Gitterpunkte in $\R^2$ mit einer guten (blauen) und einer schlechten (roten) Basis.
             Die gute Basis ist nahezu orthogonal; die schlechte Basis hat hohe Spektralnorm
             durch starke lineare Abhängigkeit.}
    \label{fig:lattice_basis}
\end{figure}

\begin{lemma}[Hadamard-Schranke]
\label{lem:hadamard}
Für jede Basis $\B = [\mathbf{b}_1, \ldots, \mathbf{b}_n]$ gilt
\[
    |\det \B| \leq \prod_{i=1}^{n} \norm{\mathbf{b}_i}.
\]
Gleichheit trifft genau dann zu, wenn die $\mathbf{b}_i$ paarweise orthogonal sind.
\end{lemma}

\begin{proof}
Dies ist die klassische Hadamard-Ungleichung, die aus der QR-Zerlegung
$\B = \mathbf{Q} \mathbf{R}$ folgt.
Da $\mathbf{Q}$ orthonormal ist, gilt $|\det \B| = |\det \mathbf{R}| = \prod_i |r_{ii}|$.
Nach der Gram-Schmidt-Prozedur ist $|r_{ii}| = \norm{\tilde{\mathbf{b}}_i} \leq \norm{\mathbf{b}_i}$,
da die Gram-Schmidt-Vektoren durch Projektion kürzer werden.
\end{proof}

\subsection{Der erste sukzessive Minimum}

\begin{definition}[Sukzessive Minima]
Für ein Gitter $\Lat \subset \R^n$ ist das $i$-te sukzessive Minimum
\[
    \lambda_i(\Lat) = \min \left\{ r \geq 0 \;\middle|\; \dim\left(\mathrm{span}\left(\Lat \cap r\mathcal{B}_n\right)\right) \geq i \right\},
\]
wobei $\mathcal{B}_n$ die $n$-dimensionale Einheitskugel bezeichnet.
Insbesondere ist $\lambda_1(\Lat)$ die \emph{Länge des kürzesten Gittervektors}.
\end{definition}

\begin{satz}[Minkowski's Erster Satz]
\label{satz:minkowski}
Für jedes volle Gitter $\Lat \subset \R^n$ gilt
\[
    \lambda_1(\Lat) \leq \sqrt{n} \cdot (\det \Lat)^{1/n}.
\]
\end{satz}

\begin{proof}
Sei $K = \lambda_1(\Lat) \cdot \mathcal{B}_n$ die abgeschlossene Kugel mit Radius $\lambda_1$.
Nehmen wir an, $K$ enthält keinen von Null verschiedenen Gitterpunkt.
Dann sind die verschobenen Mengen $\{\mathbf{x} + \tfrac{1}{2}K \mid \mathbf{x} \in \Lat\}$
paarweise disjunkt (da $\tfrac{1}{2}K$ konvex und symmetrisch).
Ihr Volumen pro Fundamentalparallelotop beträgt:
\[
    \frac{\mathrm{Vol}(\tfrac{1}{2}K)}{\det \Lat} = \frac{(\lambda_1/2)^n \cdot v_n}{\det \Lat},
\]
wobei $v_n = \pi^{n/2}/\Gamma(n/2+1)$ das Volumen der $n$-dimensionalen Einheitskugel ist.
Damit diese Mengen disjunkt sein können, muss gelten
\[
    \frac{(\lambda_1/2)^n v_n}{\det \Lat} \leq 1
    \quad \Longrightarrow \quad
    \lambda_1^n \leq \frac{2^n \det \Lat}{v_n}.
\]
Da $v_n \geq (2\pi e/n)^{n/2} \cdot (2\pi n)^{-1/2}$ und für praktische Zwecke
$v_n \geq (1/\sqrt{n})^n$ gilt, erhält man die behauptete Schranke
$\lambda_1 \leq \sqrt{n} \cdot (\det \Lat)^{1/n}$.
\end{proof}

% ============================================================
\section{Das Shortest-Vector-Problem (SVP)}
% ============================================================

\subsection{Problemdefinition}

\begin{definition}[Shortest Vector Problem -- SVP]
Gegeben eine Gitterbasis $\B \in \Z^{m \times n}$, finde einen von Null verschiedenen
Gittervektor minimaler euklidischer Länge:
\[
    \mathrm{SVP:} \quad \text{Finde } \mathbf{v} \in \Lat(\B) \setminus \{\mathbf{0}\}
                        \text{ mit } \norm{\mathbf{v}} = \lambda_1(\Lat(\B)).
\]
\end{definition}

\begin{definition}[$\gamma$-approximatives SVP ($\gamma$-SVP)]
Für einen Approximationsfaktor $\gamma \geq 1$: Finde $\mathbf{v} \in \Lat \setminus \{\mathbf{0}\}$
mit $\norm{\mathbf{v}} \leq \gamma \cdot \lambda_1(\Lat)$.
\end{definition}

\begin{figure}[ht]
    \centering
    \includegraphics[width=0.72\textwidth]{plots/plot2_svp.pdf}
    \caption{Illustration des SVP in $\R^2$: Der rote Pfeil zeigt den kürzesten Gittervektor
             $\mathbf{v}$ mit $\norm{\mathbf{v}} = \lambda_1$; der gestrichelte Kreis markiert
             den Suchradius $\lambda_1$.}
    \label{fig:svp}
\end{figure}

\subsection{Komplexität des SVP}

\begin{satz}[Ajtai 1998 -- NP-Härte von SVP]
\label{satz:svp_np}
Das exakte SVP (in der $\ell_2$-Norm) ist $\mathrm{NP}$-schwer unter Wahrscheinlichkeitsreduktionen.
\end{satz}

\begin{proof}[Beweisidee]
Ajtai~\cite{Ajtai1998} zeigte eine Worst-Case-zu-Average-Case-Reduktion:
Wenn es einen effizienten randomisierten Algorithmus für SVP im Durchschnitt gibt,
existiert auch ein Algorithmus für SVP im schlechtesten Fall.
Die Idee besteht darin, eine Zufallsinstanz von SVP zu konstruieren, die mit hoher
Wahrscheinlichkeit so schwer ist wie der schlimmstmögliche Fall.
Formal wird eine Reduktion von 3-SAT (oder einem anderen NP-vollständigen Problem)
über eine Kette von Gitterproblemen konstruiert.
Der vollständige Beweis verwendet \emph{random integer lattices} und die Tatsache,
dass das Bestimmen von $\lambda_1$ äquivalent zu Problemen in der Geometrie der Zahlen ist.
\end{proof}

\begin{satz}[Untere Schranken für SVP-Algorithmen]
\label{satz:svp_lb}
Unter der Exponential Time Hypothesis~(ETH) benötigt jeder Algorithmus für das exakte SVP
in Dimension $n$ mindestens $2^{\Omega(n)}$ Zeit.
\end{satz}

\begin{proof}
Aggarwal et al.~\cite{Aggarwal2015} bewiesen, dass das exakte SVP
Zeit $2^{n + o(n)}$ erfordert, sofern ETH gilt.
Der Beweis läuft über eine Reduktion von $k$-Sum auf SVP in passender Dimension.
Da $k$-Sum unter ETH keine subexponentielle Lösung hat, überträgt sich diese Schranke.
\end{proof}

\subsection{Algorithmen für SVP}

\subsubsection{LLL-Basisreduktion}

\begin{definition}[LLL-reduzierte Basis]
Eine Basis $(\mathbf{b}_1, \ldots, \mathbf{b}_n)$ mit Gram-Schmidt-Vektoren
$(\tilde{\mathbf{b}}_1, \ldots, \tilde{\mathbf{b}}_n)$ und Koeffizienten
$\mu_{ij} = \innerprod{\mathbf{b}_i}{\tilde{\mathbf{b}}_j}/\norm{\tilde{\mathbf{b}}_j}^2$
heißt \emph{LLL-reduziert} (mit Parameter $\delta \in (1/4, 1]$), wenn gilt:
\begin{enumerate}
    \item $|\mu_{ij}| \leq 1/2$ für alle $1 \leq j < i \leq n$ \quad (Größenreduziertheit),
    \item $\delta \norm{\tilde{\mathbf{b}}_i}^2 \leq \norm{\tilde{\mathbf{b}}_{i+1} + \mu_{i+1,i} \tilde{\mathbf{b}}_i}^2$
          für alle $i$ \quad (Lovász-Bedingung).
\end{enumerate}
\end{definition}

\begin{satz}[LLL-Algorithmus, Lenstra-Lenstra-Lovász 1982]
\label{satz:lll}
Der LLL-Algorithmus findet in Zeit $O(n^5 \log^3 \norm{\B})$
eine LLL-reduzierte Basis mit $\delta = 3/4$, bei der der erste Basisvektor
$\mathbf{b}_1$ eine $2^{(n-1)/2}$-Approximation an $\lambda_1$ liefert:
\[
    \norm{\mathbf{b}_1} \leq 2^{(n-1)/2} \cdot \lambda_1(\Lat(\B)).
\]
\end{satz}

\begin{proof}
Sei $\phi(\B) = \prod_{i=1}^{n} \norm{\tilde{\mathbf{b}}_i}^{2(n-i+1)}$ das
\emph{Potential} der Basis.
Jeder LLL-Swap reduziert $\phi$ um einen Faktor von mindestens $\delta = 3/4 > 1/2$.
Da $\phi$ ganzzahlig beschränkt ist (nach Gram-Schmidt-Argumentation)
und durch die anfängliche Größenreduktion nach oben durch $\norm{\B}_F^{n(n+1)/2}$ beschränkt,
terminiert der Algorithmus nach $O(n^2 \log \norm{\B})$ Swaps.

Für die Qualitätsschranke: In einer LLL-reduzierten Basis gilt
$\norm{\tilde{\mathbf{b}}_i}^2 \leq (4/3)^{i-1} \norm{\tilde{\mathbf{b}}_1}^2 \cdot \ldots$
Nach Ausrollen über alle $i$ ergibt sich $\norm{\mathbf{b}_1} \leq 2^{(n-1)/2} \lambda_1$.
\end{proof}

\subsubsection{Enumeration und Siebverfahren}

\begin{satz}[Kannan-Enumeration]
Die Kannan-Enumeration~\cite{Kannan1983} löst SVP exakt in Zeit
$n^{O(n)}$ und Platz $2^{O(n)}$.
\end{satz}

\begin{satz}[Siebverfahren -- AKS 2002]
Das Algorithmic-Sieve-Verfahren von Ajtai, Kumar und Sivakumar~\cite{AKS2002}
löst SVP exakt in Zeit und Platz $2^{O(n)}$.
Spätere Verfeinerungen~\cite{BDGL2016} erreichen $2^{0.292n + o(n)}$.
\end{satz}

\begin{figure}[ht]
    \centering
    \includegraphics[width=0.75\textwidth]{plots/plot4_spectral_norm.pdf}
    \caption{Spektralnorm $s_1(\B)$ guter (nahezu orthogonaler) und schlechter (hochkorrelierter)
             Basen über die Gitterdimension $n$. Die grün schraffierte Fläche zeigt die Differenz;
             ein großer Abstand erschwert Gitterproblem-Berechnungen erheblich.}
    \label{fig:spectral_norm}
\end{figure}

\begin{figure}[ht]
    \centering
    \includegraphics[width=0.75\textwidth]{plots/plot7_complexity.pdf}
    \caption{Algorithmische Komplexität für SVP/CVP-Algorithmen:
             exponentiell ($2^n$), quasi-polynomiell ($2^{(\log n)^2}$)
             und polynomiell ($n^3$).
             Für das exakte SVP sind aktuell nur (sub-)exponentielle Algorithmen bekannt.}
    \label{fig:complexity}
\end{figure}

% ============================================================
\section{GKP-Codes: Gitter im Phasenraum}
% ============================================================

\subsection{Quantenfehlerkorrektur -- Grundlagen}

Quantensysteme interagieren unvermeidlich mit ihrer Umgebung.
Ein $[[n,k,d]]$-Quantenfehlerkorrekturcode kodiert $k$ logische Qubits
in $n$ physikalischen Qubits mit Codeabstand $d$, wodurch bis zu
$\lfloor (d-1)/2 \rfloor$ Fehler korrigierbar sind~\cite{NielsenChuang2000}.

GKP-Codes arbeiten in einem \emph{kontinuierlichen Variablensystem} (Bosonenmode),
nicht mit diskreten Qubits~\cite{GKP2001}.
Die Kodierung erfolgt in $n$ bosonischen Moden (Quantenoszillatoren).

\subsection{GKP-Codeworte}

\begin{definition}[GKP-Codewörter -- quadratischer Code]
Für den einfachen quadratischen GKP-Code lauten die idealisierten Codewörter
(als Superposition von Dirac-Kämmen in der Quadraturdarstellung):
\begin{align}
    |0\rangle_{\mathrm{GKP}} &= \sum_{n \in \Z} |2n\sqrt{\pi}\rangle_q,  \label{eq:gkp0}\\
    |1\rangle_{\mathrm{GKP}} &= \sum_{n \in \Z} |(2n+1)\sqrt{\pi}\rangle_q. \label{eq:gkp1}
\end{align}
In der Praxis werden diese durch eine Gaußsche Hülle $e^{-q^2/(2\sigma^2)}$
regularisiert~\cite{GKP2001}.
\end{definition}

\begin{figure}[ht]
    \centering
    \includegraphics[width=0.80\textwidth]{plots/plot3_gkp_codewords.pdf}
    \caption{GKP-Codewörter $|0\rangle$ (blau) und $|1\rangle$ (rot) im Quadraturraum $q$,
             dargestellt als endliche Überlagerung Gaußscher Pakete ($\sigma = 2$, $|n| \leq 10$).
             Die schwarze gestrichelte Kurve zeigt die Gaußsche Gesamthülle.}
    \label{fig:gkp_codewords}
\end{figure}

\subsection{Gittersicht auf GKP-Codes}

\subsubsection{Verschiebungsoperator und Stabilisatorgruppe}

\begin{definition}[Verschiebungsoperator]
Für $\xi = (\xi_q, \xi_p)^T \in \R^{2n}$ ist der Verschiebungsoperator
\[
    D(\xi) = \exp\!\left\{-i\sqrt{2\pi}\,\xi^T J \hat{x}\right\},
    \quad J = \begin{pmatrix} 0 & I_n \\ -I_n & 0 \end{pmatrix},
\]
wobei $\hat{x} = (\hat{q}_1, \ldots, \hat{q}_n, \hat{p}_1, \ldots, \hat{p}_n)^T$
der Quadraturvektor und $J$ die Standard-Symplektikmatrix ist.
\end{definition}

\begin{definition}[Stabilisatorgitter]
Die Stabilisatorgruppe des GKP-Codes ist
\[
    \mathcal{S} = \{D(\xi_1), \ldots, D(\xi_{2n})\}.
\]
Diese ist isomorph zu einem Gitter
\[
    \Lat = \{\xi \in \R^{2n} \mid \xi = a^T M,\; a \in \Z^{2n \times 2n}\},
\]
wobei $M$ die Generatormatrix des Gitters ist.
\end{definition}

\begin{definition}[Symplektisches Dualgitter]
Das \emph{symplektische Dualgitter} von $\Lat$ ist
\[
    \Lat^\perp = \left\{ \xi^\perp \in \R^{2n} \;\middle|\;
                         (\xi^\perp)^T J \xi \in \Z \;\;\forall \xi \in \Lat \right\}.
\]
\end{definition}

\begin{figure}[ht]
    \centering
    \includegraphics[width=0.72\textwidth]{plots/plot8_symplectic_dual.pdf}
    \caption{Primärgitter $\Lat$ (blaue Punkte/Pfeile) und symplektisches Dualgitter $\Lat^\perp$
             (rote Punkte/Pfeile) in $\R^2$.
             Das Dualgitter ist dichter; der Quotient $\Lat^\perp / \Lat$ parametrisiert
             die logisch distinkten Verschiebungen (Fehlerklassen).}
    \label{fig:symplectic_dual}
\end{figure}

\begin{lemma}[Charakterisierung des Coderaums]
\label{lem:coderaum}
Der Codezustand $|\psi\rangle$ liegt genau dann im GKP-Coderaum, wenn
$D(\xi)|\psi\rangle = |\psi\rangle$ für alle $\xi \in \Lat$ gilt,
d.h.\ alle Stabilisatoren trivial auf $|\psi\rangle$ wirken.
\end{lemma}

\begin{proof}
Da die Stabilisatoren $D(\xi)$ für $\xi \in \Lat$ eine abelsche Gruppe bilden
(kommutativität folgt aus $[D(\xi), D(\xi')] = e^{i\omega(\xi,\xi')}$ mit
$\omega(\xi,\xi') = \xi^T J \xi' \in \Z$ für $\xi, \xi' \in \Lat$),
definiert ihr gemeinsamer $+1$-Eigenraum den Codezustand.
Die Abgeschlossenheit unter Multiplikation garantiert, dass dieser Eigenraum
nicht-trivial ist und genau die gewünschten logischen Qubits trägt.
\end{proof}

\subsection{GKP-Codeabstand}

\begin{definition}[GKP-Codeabstand]
Der \emph{Codeabstand} des GKP-Codes ist
\[
    \Delta = \Delta(\Lat) := \min_{0 \neq x \in \Lat^\perp / \Lat} \norm{x}.
\]
\end{definition}

\begin{satz}[Eigenschaften guter GKP-Codes]
\label{satz:good_gkp}
Ein GKP-Code heißt \emph{gut}, wenn
\begin{enumerate}
    \item $\lim_{n \to \infty} \frac{\log \det(\mathcal{L}_n)}{n} > 0$, und
    \item $\Delta^2 = \Omega(n)$.
\end{enumerate}
\end{satz}

\begin{proof}
\textbf{Bedingung~1:} Der Nenner $n$ skaliert mit der Systemgröße.
Wenn $\det \Lat_n \geq c^n$ für eine Konstante $c > 1$, wächst
$\log \det \Lat_n \geq n \log c > 0$ linear.
Dies entspricht einer konstanten \emph{Ratendichte} und ist das Analogon zur positiven
Coderate bei klassischen Codes.

\textbf{Bedingung~2:}
Für $\Delta^2 = \Omega(n)$ muss der kürzeste Vektor in $\Lat^\perp / \Lat$
eine Länge $\Omega(\sqrt{n})$ haben.
Nach Satz~\ref{satz:minkowski} gilt $\lambda_1(\Lat^\perp / \Lat) \lesssim \sqrt{n}$;
ein guter Code schöpft diese Schranke also bis auf konstante Faktoren aus.
\end{proof}

\subsection{GKP-Codes aus kryptographischen Gittern}

\subsubsection{NTRU-Gitter und GKP}

\begin{satz}[GKP-Codes aus NTRU-Gittern, Conrad et al.~2024]
\label{satz:ntru_gkp}
Sei $H = \begin{pmatrix} I_n & A \\ 0 & qI_n \end{pmatrix}$ $q$-symplektisch, wobei $A$
symmetrisch ist. Dann erzeugt $H$ einen GKP-Code.
Mit einer zufälligen NTRU-Matrix $A_\phi(h_1^{-1}h_2)$ erreicht der resultierende Code
für $n \leq q^2/\mathrm{const}$ den Codeabstand
\[
    \Delta \propto \min\!\left\{ \sqrt{n/q},\; \sqrt{q} \right\}.
\]
\end{satz}

\begin{proof}[Beweisidee]
Die $q$-Symplektizität von $H$ garantiert $H^T J H = qJ$, was die Kommutatorbedingung
für Stabilisatoren erfüllt (Verschiebungsoperatoren, die im Gitter liegen, kommutieren).
Die Codeabstandsschranke folgt aus der typischen Länge kürzester Vektoren in
zufälligen NTRU-Gittern, die durch Number-Field-Sieve-Argumente bekannt sind.
\end{proof}

\subsubsection{Ring-SIS Gitter und GKP}

\begin{satz}[GKP-Codes aus Ring-SIS-Gittern]
\label{satz:ring_sis_gkp}
Sei $H = \begin{pmatrix} I_n & A_\phi(h_1^{-1}h_2) \\ 0 & qI_n \end{pmatrix}$
mit zufälliger Ring-SIS-Matrix.
Der resultierende GKP-Code erreicht den Codeabstand
\[
    \Delta \propto \min\!\left\{\sqrt{n},\; \sqrt{q}\right\}
\]
für $n \leq \mathrm{poly}(q)$ und definiert eine Familie guter GKP-Codes.
\end{satz}

\begin{proof}
Die Konstruktion basiert auf der Tatsache, dass Ring-SIS-Gitter strukturell
mehr Symmetrien besitzen als allgemeine SIS-Gitter.
Insbesondere gilt für zufällige $h_1, h_2$ im Ring $\Z_q[x]/(f(x))$, dass
der kürzeste Vektor des Gitters typischerweise Länge $\Omega(\sqrt{n})$ hat
(nach randomisierten Gitterargumenten aus~\cite{LyubashevskyMicciancio2006}).
Da $n \leq \mathrm{poly}(q)$, ist $\sqrt{n} \leq q^{c/2}$ für ein $c \geq 1$;
der Minimus ist daher $\sqrt{n}$.
\end{proof}

\begin{figure}[ht]
    \centering
    \includegraphics[width=0.78\textwidth]{plots/plot5_gkp_distance.pdf}
    \caption{GKP-Codeabstand $\Delta \propto \min(\sqrt{n}, \sqrt{q})$ aus der
             Ring-SIS-Konstruktion in Abhängigkeit der Dimension $n$ für verschiedene Moduli $q$.
             Die gestrichelten Vertikallinien markieren die Ãœbergangspunkte $n = q$.}
    \label{fig:gkp_distance}
\end{figure}

\subsection{Skalierte GKP-Codes}

\begin{definition}[Skalierter GKP-Code]
Sei $M_0$ eine selbst-duale Matrix mit $M_0^T J M_0 = J$.
Ein \emph{skalierter GKP-Code} verwendet die Generatormatrix
$M = \sqrt{\lambda} M_0$, sodass $M^T J M = \lambda J$.
Man nennt einen solchen Code \emph{$\lambda$-symplektisch}.
\end{definition}

\begin{lemma}[Selbst-Dualität skalierter Codes]
Für einen selbst-dualen Basiscode mit $M_0^T J M_0 = J$ gilt:
\[
    M^T J M = \lambda J \quad \Longrightarrow \quad \Delta(M) = \sqrt{\lambda} \cdot \Delta(M_0).
\]
\end{lemma}

\begin{proof}
Der Codeabstand skaliert linear mit dem Skalierungsfaktor $\sqrt{\lambda}$,
da alle Gittervektoren um $\sqrt{\lambda}$ skaliert werden.
Die Minimumbedingung $\min_{0 \neq x \in \Lat^\perp/\Lat}\norm{x}$ skaliert daher als
$\sqrt{\lambda} \cdot \min_{0 \neq x \in \Lat_0^\perp/\Lat_0}\norm{x} = \sqrt{\lambda}\,\Delta(M_0)$.
\end{proof}

% ============================================================
\section{Gitterbasiertes Anonymes Reputationssystem}
% ============================================================

\subsection{Systemmodell und Sicherheitsziele}

\begin{definition}[Anonymes Reputationssystem]
Ein \emph{anonymes Reputationssystem} ist ein kryptographisches Protokoll
mit Teilnehmern Verkäufer $S$, Nutzer $U$, Vertrauenspartei $T$ und Öffentlichkeit $P$,
das folgende Sicherheitseigenschaften erfüllt:
\begin{enumerate}
    \item \textbf{Anonymität:} Niemand außer $T$ kann bestimmen, welcher Nutzer eine Bewertung schrieb.
    \item \textbf{Autorisierung:} Nur von $S$ autorisierte Nutzer können Bewertungen schreiben.
    \item \textbf{Einfachverwendbarkeit:} Jede Autorisierung kann nur einmal zum Schreiben verwendet werden.
    \item \textbf{Rückverfolgbarkeit:} $T$ kann bei Missbrauch die Identität eines Nutzers aufdecken.
\end{enumerate}
\end{definition}

\begin{figure}[ht]
    \centering
    \includegraphics[width=0.82\textwidth]{plots/plot6_reputation_workflow.pdf}
    \caption{Protokollablauf des anonymen Reputationssystems nach~\cite{BBP23}:
             Der Verkäufer autorisiert den Nutzer; dieser schreibt anonym eine Bewertung,
             die öffentlich verifizierbar ist. Bei Missbrauch deanonymisiert die Vertrauenspartei.}
    \label{fig:reputation}
\end{figure}

\subsection{Kryptographische Bausteine}

\begin{definition}[Digitale Gittersignatur]
Eine Gittersignatur über $\Z_q^n$ ist ein Tupel $(\mathbf{A}, \mathbf{t}, \mathbf{s})$ mit:
\begin{itemize}
    \item Öffentlichem Schlüssel: $\mathbf{A} \in \Z_q^{n \times m}$, $\mathbf{t} = \mathbf{A}\mathbf{s} \pmod q$
    \item Geheimem Schlüssel: $\mathbf{s} \in \Z^m$ mit $\norm{\mathbf{s}} \leq \beta$
    \item Signatur: $\mathbf{z} \in \Z^m$ mit $\norm{\mathbf{z}} \leq B$ und $\mathbf{A}\mathbf{z} = \mathbf{t} \pmod q$
\end{itemize}
Die Sicherheit basiert auf der Schwierigkeit von SIS/LWE.
\end{definition}

\begin{definition}[Learning With Errors (LWE)]
Für Sicherheitsparameter $n$, Modulus $q$ und Fehlerparameter $\alpha$ ist LWE das Problem,
$\mathbf{s} \in \Z_q^n$ aus Stichproben $(\mathbf{a}_i, \mathbf{a}_i^T \mathbf{s} + e_i)$
mit $\mathbf{a}_i \leftarrow \Z_q^n$ gleichverteilt und $e_i \sim \mathcal{D}_{\Z,\alpha q}$
zu bestimmen.
\end{definition}

\begin{satz}[Sicherheit von LWE -- Regev 2005]
\label{satz:lwe_sicherheit}
LWE mit Parametern $(n, q, \alpha)$ ist mindestens so schwer wie das
$\tilde{O}(n/\alpha)$-approximative SVP und CVP in der schlimmsten Eingabe
(Worst-Case).
\end{satz}

\begin{proof}[Beweisidee]
Regev~\cite{Regev2005} konstruiert eine quantenbasierte Reduktion:
Gegeben ein Algorithmus $\mathcal{A}$ für LWE, baue einen Algorithmus für
$\gamma$-SVP mit $\gamma = \tilde{O}(n/\alpha)$.
Die Reduktion verläuft über eine Kette
\[
    \gamma\text{-SVP} \;\xrightarrow{\text{Gaußsches Sampling}} \;
    \text{BDD} \;\xrightarrow{\text{Quantenreduktion}} \;
    \text{LWE},
\]
wobei BDD (Bounded Distance Decoding) als Zwischenproblem fungiert.
Quantencomputer werden für die Gausssche-Sampling-Phase benötigt;
die Reduktion läuft somit auf einem Quantencomputer, beweist aber die Schwierigkeit
des klassischen LWE relativ zu SVP.
\end{proof}

\subsection{Anonyme Signatures und Tag-Schemata}

\begin{definition}[Tag-Schema]
Ein Tag-Schema $(\mathsf{Tag.Setup}, \mathsf{Tag.Sign}, \mathsf{Tag.Verify})$ ordnet
jeder Autorisierungssitzung ein eindeutiges Tag $\tau$ zu, sodass:
\begin{itemize}
    \item Zwei Bewertungen der gleichen Sitzung dasselbe Tag $\tau$ haben,
    \item das Tag $\tau$ keine Information über die Nutzeridentität preisgibt.
\end{itemize}
\end{definition}

\begin{satz}[Nicht-Bindbarkeit unter LWE]
\label{satz:non_frameability}
Unter der LWE-Annahme kann kein effizienter Angreifer
eine gültige, von $T$ akzeptierte Bewertung mit einem anderen Tag als seinem eigenen
verknüpfen (\emph{Non-Frameability}).
\end{satz}

\begin{proof}
Sei $\mathcal{A}$ ein Angreifer, der Non-Frameability bricht:
Er produziert eine gültige Bewertung $(m, \sigma, \tau)$ für ein Tag $\tau \neq \tau_{\mathcal{A}}$.
Dies impliziert, dass $\mathcal{A}$ eine LWE-Instanz löst, denn:
\begin{enumerate}
    \item Das Tag $\tau$ ist kryptographisch an die Autorisierungssitzung gebunden.
    \item Das Fälschen eines Tags zu einer fremden Sitzung erfordert das Auffinden eines
          LWE-Geheimnisses $\mathbf{s}$ ohne Kenntnis des geheimen Schlüssels $T$.
    \item Dies wäre ein Widerspruch zu Satz~\ref{satz:lwe_sicherheit}.
\end{enumerate}
\end{proof}

\subsection{Zero-Knowledge-Beweise über Gitter}

\begin{definition}[Nicht-interaktiver Zero-Knowledge-Beweis (NIZK)]
Ein NIZK für eine Aussage $x \in L$ ist ein Beweis $\pi$, sodass:
\begin{itemize}
    \item \textbf{Vollständigkeit:} Für alle $x \in L$ mit Zeugen $w$ gilt $\Pr[\mathsf{Verify}(x, \pi) = 1] = 1$.
    \item \textbf{Solidität:} Für $x \notin L$ gilt $\Pr[\mathsf{Verify}(x, \pi) = 1] \leq \epsilon_{\mathrm{sound}}$.
    \item \textbf{Zero-Knowledge:} $\pi$ gibt keinerlei Information über $w$ preis.
\end{itemize}
\end{definition}

\begin{satz}[Gitterbasierter NIZK -- Sicherheit]
\label{satz:nizk}
Das gitterbasierte anonyme Reputationssystem aus~\cite{BBP23} ist sicher gegen
klassische Angreifer unter den Annahmen LWE und SIS.
Gegen Quantenangreifer besteht eine verbleibende Unsicherheit bezüglich einer
spezifischen Technik im Sicherheitsbeweis.
\end{satz}

\begin{proof}[Beweisidee]
Die Sicherheit ergibt sich aus dem modularen Aufbau:
\begin{enumerate}
    \item Anonymität $\Leftarrow$ Verschlüsselungsschema (IND-CPA unter LWE).
    \item Autorisierung $\Leftarrow$ Digitale Signatur (EUF-CMA unter SIS).
    \item Einfachverwendbarkeit $\Leftarrow$ Tag-Schema (Bindbarkeit unter LWE).
    \item Rückverfolgbarkeit $\Leftarrow$ Vertrauensparteischlüssel (Revokation über Gitterbasis).
\end{enumerate}
Die verbleibende Quantenunsicherheit betrifft den NIZK-Simulator,
der aktuell einen Quantenangreifer nicht vollständig simulieren kann.
\end{proof}

% ============================================================
\section{Verbindung zum Subgraph Algorithmus}
% ============================================================

\subsection{Gitterstrukturen als Graphen}

\begin{definition}[Gitter-Cayley-Graph]
Für ein Gitter $\Lat(\B)$ mit Basis $\B$ ist der \emph{Cayley-Graph}
$G_\Lat = (V_\Lat, E_\Lat)$ definiert durch:
\begin{itemize}
    \item $V_\Lat = \Lat(\B)$ (Gitterpunkte als Knoten),
    \item $(u, v) \in E_\Lat \iff v - u \in \{\pm \mathbf{b}_1, \ldots, \pm \mathbf{b}_n\}$ (benachbart nach Basisvektoren).
\end{itemize}
\end{definition}

\begin{satz}[Subgraph Algorithmus auf Gitter-Cayley-Graphen]
\label{satz:subgraph_gitter}
Sei $\Lat_1 \subseteq \Lat_2$ (Untergitter). Dann gilt $G_{\Lat_1} \subseteq G_{\Lat_2}$
als Subgraph, und der Subgraph Algorithmus~\cite{Epp2026Subgraph} erkennt diese
Beziehung in Zeit $O(n^3)$.
\end{satz}

\begin{proof}
Eine Untergitter-Beziehung $\Lat_1 \subseteq \Lat_2$ bedeutet, dass jede Basis von $\Lat_1$
durch ganzzahlige Linearkombinationen aus der Basis von $\Lat_2$ darstellbar ist.
Die Adjazenzmatrizen $A_1$ und $A_2$ der entsprechenden Cayley-Graphen erfüllen dann
$A_1 \subseteq A_2$ (komponentenweise).
Der Subgraph Algorithmus berechnet für beide Matrizen Signatur-Arrays und vergleicht diese
mittels LCS-basierter Rotation, was in $O(n^3)$ ausführbar ist.
Die Korrektheit folgt aus Lemma~1 in~\cite{Epp2026Subgraph} (Injektivität der Signaturen).
\end{proof}

\begin{korollar}
Die Basisreduktionsproblematik (Ãœbergang von schlechter zu guter Basis)
entspricht dem Suchen eines isomorphen Subgraphen im Cayley-Graphen.
\end{korollar}

\begin{bemerkung}
Für den GKP-Code gilt: Die symplektische Dualgitter-Relation $\Lat \subseteq \Lat^\perp$
ist direkt als Subgraph-Relation der entsprechenden Cayley-Graphen erkennbar.
Der Subgraph Algorithmus kann damit zur effizienten Überprüfung von GKP-Codebedingungen
eingesetzt werden.
\end{bemerkung}

% ============================================================
\section{Zusammenfassung}
% ============================================================

Diese Arbeit hat drei zentrale Resultate aus Gittertheorie, Quantenfehlerkorrektur
und gitterbasierter Kryptographie formal erarbeitet:

\begin{enumerate}
    \item \textbf{SVP:}
          Das Shortest-Vector-Problem ist NP-schwer (Satz~\ref{satz:svp_np})
          und benötigt unter ETH $2^{\Omega(n)}$ Zeit (Satz~\ref{satz:svp_lb}).
          Der LLL-Algorithmus liefert eine $2^{(n-1)/2}$-Approximation in Polynomialzeit
          (Satz~\ref{satz:lll}).

    \item \textbf{GKP-Codes:}
          GKP-Codewörter sind über Gitter im symplektischen Phasenraum definiert.
          Ring-SIS-Konstruktionen erreichen Codeabstand $\Delta \propto \min(\sqrt{n}, \sqrt{q})$
          (Satz~\ref{satz:ring_sis_gkp}), was eine Familie guter Codes definiert.

    \item \textbf{Anonyme Reputationssysteme:}
          Unter LWE/SIS-Annahmen erreicht das System aus~\cite{BBP23} alle vier
          Sicherheitsziele gegen klassische Angreifer
          (Sätze~\ref{satz:non_frameability}, \ref{satz:nizk}).
\end{enumerate}

% ============================================================
\section{Ausblick}
% ============================================================

Die in dieser Arbeit behandelten Themengebiete -- Gitterprobleme, GKP-Codes und
gitterbasierte Kryptographie -- sind Gegenstand intensiver aktueller Forschung.
Im Folgenden werden die wichtigsten offenen Fragen und zukünftigen Forschungsrichtungen
ausführlich diskutiert.

\subsection{SVP und Gitterbasisreduktion}

\subsubsection{Verbesserte Approximationsalgorithmen}

Der LLL-Algorithmus liefert eine $2^{O(n)}$-Approximation in Polynomialzeit.
Eine fundamentale offene Frage ist, ob eine Approximation mit Faktor $n^{O(1)}$
(d.h.\ polynomiell in $n$) in Polynomialzeit möglich ist.
Das BKZ-Verfahren (Block Korkine-Zolotarev)~\cite{Schnorr1994} interpoliert
zwischen LLL und exaktem SVP durch einen Blockparameter $\beta$:
Größeres $\beta$ liefert bessere Approximationen auf Kosten exponentieller
Laufzeit in $\beta$.
Die optimale Wahl von $\beta$ in Abhängigkeit von Dimension und Sicherheitsanforderung
ist weiterhin aktives Forschungsgebiet~\cite{Chen2011}.

Ein vielversprechender Ansatz ist die Verwendung von \emph{maschinellem Lernen}
zur Steuerung des BKZ-Algorithmus: neuronale Netze können lernen, welche
Blockpositionen besonders fruchtbar für Reduktionsschritte sind~\cite{Novaes2023}.

\subsubsection{Quantenalgorithmen für SVP}

Shor's Algorithmus löst diskrete Logarithmen und Faktorisierung,
greift jedoch SVP nicht an.
Der Quantenalgorithmus von Hallgren~\cite{Hallgren2007} löst Pell's Gleichung
und verwandte Probleme auf idealen Gittern, ist aber nicht auf allgemeine Gitter übertragbar.
Für allgemeine Gitter gibt es bislang keinen subexponentiellen Quantenalgorithmus.
Die Klärung der Frage, ob Quantencomputer SVP deutlich beschleunigen können,
wäre ein Durchbruch sowohl für die Gitterkryptographie als auch die Quantenkomplexitätstheorie.

\subsubsection{Härtere Sicherheitsreduktionen}

Die Lücke zwischen der NP-Härte für exaktes SVP und den polynomiellen Approximationsfaktoren,
auf denen kryptographische Konstruktionen beruhen (typischerweise $\gamma = n^{O(1)}$),
ist bislang nicht geschlossen.
Es ist offen, ob das $\sqrt{n}$-approximative SVP ebenfalls NP-schwer ist~\cite{Khot2005}.

\subsection{GKP-Codes und Quantenfehlerkorrektur}

\subsubsection{Vorbereitung von GKP-Zuständen}

Die idealisierten GKP-Codewörter aus \eqref{eq:gkp0} und \eqref{eq:gkp1}
sind nicht normierbar; in der Praxis müssen sie durch endliche
Gaußsche Überlagerungen approximiert werden.
Die Vorbereitung dieser Zustände mit ausreichender Güte
(gemessen am Quetschen) bleibt eine große experimentelle Herausforderung~\cite{Campagne2020}.
Aktuelle Experimente erreichen ein Quetschen von ca.~10--12\,dB in supraleitenden
Systemen und photonischen Plattformen.
Für fehlertolerante Quantenberechnungen werden schätzungsweise 20--30\,dB benötigt.

\subsubsection{Erweiterung auf höherdimensionale GKP-Codes}

Der in dieser Arbeit betrachtete quadratische GKP-Code in $n$ Moden kodiert $n$ logische Qubits.
Es ist bekannt~\cite{Harrington2001}, dass höherdimensionale Gitter (z.B.\ hexagonale oder
$E_8$-Gitter) bessere Fehlerkorrekturkapazitäten bieten.
Die systematische Konstruktion optimaler GKP-Codes aus zahlentheoretischen Gittern
(z.B.\ Barnes-Wall-Gitter, Leech-Gitter) ist ein aktives Forschungsfeld~\cite{Royer2022}.

\subsubsection{GKP-Codes und Fehlerkorrektur-Schwellenwerte}

Für fault-tolerante Quantencomputer sind Fehlerschwellen entscheidend:
Wenn die physikalische Fehlerrate unter dem Schwellenwert liegt, kann beliebig
lange gerechnet werden.
GKP-Codes haben nachweislich hohe Schwellenwerte~\cite{Menicucci2014},
insbesondere wenn sie mit anderen Codes wie Surface-Codes kombiniert werden.
Die genaue Quantifizierung dieser Schwellenwerte für realistische Fehlermodelle
(Quetsch-Rauschen, Verlust, thermisches Rauschen) bleibt Gegenstand aktiver Forschung.

\subsubsection{Verbindung zwischen GKP-Codeabstand und Gitterparametern}

Die Verbindung zwischen dem GKP-Codeabstand $\Delta$ und dem ersten sukzessiven Minimum
$\lambda_1$ des Dualgitters ist formal etabliert.
Die optimale Wahl von Gitterparametern (Determinante, Kissing-Zahl, Packunsdichte)
zur Maximierung von $\Delta$ ist jedoch weiterhin offen und berührt klassische Probleme
der Gitterpackungstheorie~\cite{ConwaySloane1993}.

\subsection{Gitterbasierte Kryptographie}

\subsubsection{NIST-Standardisierung und praktische Sicherheitsparameter}

NIST hat 2024 die ersten post-quanten Kryptographiestandards veröffentlicht:
CRYSTALS-Kyber (ML-KEM), CRYSTALS-Dilithium (ML-DSA) und FALCON (FN-DSA)
basieren alle auf dem Module-LWE bzw.\ NTRU-Problem~\cite{NIST2024}.
Die Wahl konkreter Sicherheitsparameter (Dimension $n$, Modulus $q$, Rauschparameter $\sigma$)
ist ein Abwägen zwischen Sicherheit und Effizienz.
Aufbauend auf der NIST-Standardisierung werden weiterhin Implementierungen
auf eingebetteten Systemen (Stichwort: AUTOSAR, ARM Cortex)
und die Abschätzung konkreter Angriffslaufzeiten~\cite{Bai2021} diskutiert.

\subsubsection{Vollhomomorphe Verschlüsselung (FHE)}

FHE ermöglicht Berechnungen auf verschlüsselten Daten.
Das CKKS-Schema~\cite{CKKS2017} und BFV/BGV~\cite{BFV2012} basieren alle auf LWE/RLWE.
Aktuelle Forschungsrichtungen umfassen:
\begin{itemize}
    \item Effiziente Bootstrapping-Verfahren zur unbeschränkten Rechentiefe.
    \item Homomorphe Auswertung neuronaler Netze.
    \item FHE auf GKP-Basis mit kontinuierlichen Variablen.
\end{itemize}

\subsubsection{Gruppen-Signaturen und anonyme Credentials}

Das in dieser Arbeit besprochene anonyme Reputationssystem ist ein Spezialfall
eines allgemeinen gitterbasierten Gruppen-Signatur-Schemas.
Aktuelle Forschung zielt auf kompaktere Beweise, effizientere Revokationsmechanismen
und Erweiterungen für mehrere Vertrauensparteien (Threshold-Deanonymisierung)~\cite{Lyubashevsky2022}.
Die Quantensicherheit der NIZK-Komponente (Fiat-Shamir-Transformation in der
Random-Oracle-Modell-Variante) bleibt ein wichtiges offenes Problem.

\subsubsection{Gitter und blockchain-basierte Systeme}

Blockchain-Systeme nutzen derzeit elliptische Kurvenkryptographie.
Die Migration zu gitterbasierten Signaturen (z.B.\ FALCON oder Dilithium) ist
aus Quantensicherheitsgründen notwendig.
Besondere Herausforderungen sind:
\begin{itemize}
    \item Größere Signaturgrößen (Dilithium: $\sim$2\,kB vs.\ ECDSA: 64\,Bytes).
    \item Kompatibilität mit bestehenden Smart-Contract-Plattformen.
    \item Effizienz auf ressourcenbeschränkten IoT-Knoten.
\end{itemize}

\subsection{Algorithmische Aspekte und der Subgraph Algorithmus}

\subsubsection{Parallelisierung des Subgraph Algorithmus für Gitter}

Der Subgraph Algorithmus~\cite{Epp2026Subgraph} hat eine Laufzeit $O(n^3)$,
die durch Parallelisierung auf $O(n^2)$ pro Prozessor reduziert werden kann
(jede Rotation unabhängig berechnet).
Für Gitter mit $n > 10\,000$ (typisch in post-quanten Kryptographie)
wären GPU-Implementierungen oder spezialisierte Hardware~(FPGA/ASIC) notwendig.

\subsubsection{Approximative Subgraph-Suche in Gittern}

Für SVP-Approximation ist es nicht notwendig, exakte Subgraph-Isomorphismen zu finden.
Ein approximativer Subgraph Algorithmus, der eine $\gamma$-Ähnlichkeit erkennt,
könnte als Basis für neuartige Approximationsalgorithmen für SVP dienen.
Formale Verbindungen zwischen der LCS-Schranke des Subgraph Algorithmus
und dem Approximationsfaktor von Gitteralgorithmen sind zu untersuchen.

\subsubsection{Subgraph Algorithmus und Graph-Neuronale-Netze}

Die Signatur-Array-Darstellung von Graphen im Subgraph Algorithmus
lässt sich als Feature-Vektor für Graph-Neuronale-Netze (GNN) interpretieren~\cite{Xu2019}.
Eine Verbindung zwischen der LCS-basierten Ähnlichkeit im Subgraph Algorithmus
und Aggregationsfunktionen in GNNs könnte neue Perspektiven auf
Gitter-Klassifikationsprobleme eröffnen.

\subsection{Experimentelle und ingenieurwissenschaftliche Aspekte}

\subsubsection{Implementierung auf eingebetteten Systemen}

Die Implementierung gitterbasierter Kryptographie auf eingebetteten Systemen
(ESP32, ARM Cortex-M) ist aufgrund begrenzten Speichers und Rechenleistung
herausfordernd~\cite{Epp2026Embedded}.
Besonders relevant sind:
\begin{itemize}
    \item Optimierung der NTT (Number Theoretic Transform) für modular arithmetic.
    \item Seitenkanalangriff-resistente Implementierungen (Power Analysis).
    \item Integration in AUTOSAR-konforme Software-Architekturen.
\end{itemize}

\subsubsection{GKP-Codes auf Photonenplattformen}

Das Paderborner PhoQS-Institut (Photonic Quantum Systems) arbeitet an der
experimentellen Realisierung von GKP-Codes auf linearen optischen Systemen.
Offene Herausforderungen sind:
\begin{itemize}
    \item Messung der Quadraturen mit ausreichender Präzision.
    \item Echtzeit-Dekodierung mit begrenzter Klassik-Rechenzeit.
    \item Skalierung auf $>10$ Moden für mehrdimensionale GKP-Codes.
\end{itemize}

\subsection{Verbindungen zu weiteren Forschungsfeldern}

\subsubsection{Gitter und Zahlentheorie}

Viele Gitterprobleme (SVP auf idealen Gittern, Short Integer Solution in Polynomringen)
haben direkte Verbindungen zur algebraischen Zahlentheorie~\cite{PeikertRosen2007}.
Die Nutzung von Euklidischen Zahlkörpern ($\mathbb{Q}(\zeta_n)$) ermöglicht
besonders effiziente Implementierungen durch NTT-freundliche Primzahlen.

\subsubsection{Kodierungstheorie und Gitter}

GKP-Codes verknüpfen Gittertheorie und Quantenfehlerkorrektur direkt.
Die analoge Frage für klassische Codes -- welche Gitter entsprechen guten
klassischen linearen Codes? -- führt zur Konstruktion von LDPC-Codes aus Gittern~\cite{Sipser1996}.
Diese Verbindung ist sowohl für klassische als auch quantenkohärente Systeme fruchtbar.

\subsubsection{Kryptographie und Physik}

Die Verbindung zwischen physikalischen Symmetrien (symplektische Gruppe $\mathrm{Sp}(2n,\R)$,
die Quantenmechanik des harmonischen Oszillators) und kryptographischen Strukturen
(Gitter im Phasenraum, GKP-Codes) ist fundamental.
Tiefere Einsichten in die Verbindung zwischen Entropie, Quetschen und
kryptographischer Sicherheit könnten neue Angriffsvektoren oder Beweistechniken liefern.

% ============================================================
\section*{Danksagung}
% ============================================================

Der Autor dankt dem PhoQS-Institut der Universität Paderborn für die
Bereitstellung des Poster-Materials zu GKP-Codes und anonymen Reputationssystemen,
das als Ausgangspunkt dieser Arbeit diente.

% ============================================================
\begin{thebibliography}{99}
% ============================================================

\bibitem{Ajtai1996}
M.~Ajtai.
\newblock Generating hard instances of lattice problems.
\newblock In \textit{Proc. STOC}, pages 99--108, 1996.

\bibitem{Ajtai1998}
M.~Ajtai.
\newblock The shortest vector problem in $L_2$ is NP-hard for randomized reductions.
\newblock In \textit{Proc. STOC}, pages 10--19, 1998.

\bibitem{Aggarwal2015}
D.~Aggarwal, D.~Dadush, O.~Regev, and N.~Stephens-Davidson.
\newblock Solving the Shortest Vector Problem in $2^n$ Time via Discrete
  Gaussian Sampling.
\newblock In \textit{Proc. STOC}, pages 733--742, 2015.

\bibitem{AKS2002}
M.~Ajtai, R.~Kumar, and D.~Sivakumar.
\newblock A sieve algorithm for the shortest lattice vector problem.
\newblock In \textit{Proc. STOC}, pages 601--610, 2002.

\bibitem{Bai2021}
S.~Bai et al.
\newblock Concrete security of CRYSTALS-Kyber.
\newblock \textit{IACR ePrint}, 2021/1058.

\bibitem{BBP23}
J.~Blömer, J.~Bobolz, and L.~Porzenheim.
\newblock A generic construction of an anonymous reputation system and
  instantiations from lattices.
\newblock In \textit{ASIACRYPT}, Springer, 2023.
\newblock \url{https://ia.cr/2023/464}.

\bibitem{BDGL2016}
A.~Becker, L.~Ducas, N.~Gama, and T.~Laarhoven.
\newblock New directions in nearest neighbor searching with applications to
  lattice sieving.
\newblock In \textit{Proc. SODA}, pages 10--24, 2016.

\bibitem{BFV2012}
F.~Brakerski, C.~Gentry, and V.~Vaikuntanathan.
\newblock (Leveled) Fully Homomorphic Encryption without Bootstrapping.
\newblock In \textit{Proc. ITCS}, 2012.

\bibitem{Campagne2020}
R.~Campagne-Ibarcq et al.
\newblock Quantum error correction of a qubit encoded in grid states of an
  oscillator.
\newblock \textit{Nature}, 584:368--372, 2020.

\bibitem{Chen2011}
Y.~Chen and P.~Q.~Nguyen.
\newblock BKZ 2.0: Better Lattice Security Estimates.
\newblock In \textit{ASIACRYPT}, Springer, 2011.

\bibitem{CKKS2017}
J.~H.~Cheon, A.~Kim, M.~Kim, and Y.~Song.
\newblock Homomorphic Encryption for Arithmetic of Approximate Numbers.
\newblock In \textit{ASIACRYPT}, 2017.

\bibitem{ConwaySloane1993}
J.~H.~Conway and N.~J.~A.~Sloane.
\newblock \textit{Sphere Packings, Lattices and Groups}.
\newblock Springer, 3rd edition, 1993.

\bibitem{Epp2026Subgraph}
S.~Epp.
\newblock Der Subgraph Algorithmus.
\newblock, Januar 2026.
\newblock \url{https://github.com/hjstephan86/science}.

\bibitem{Epp2026Embedded}
S.~Epp.
\newblock Design Optimization in Embedded Systems via Subgraph Isomorphism
  and Pareto Optimization.
\newblock, 2026.

\bibitem{GKP2001}
D.~Gottesman, A.~Kitaev, and J.~Preskill.
\newblock Encoding a qubit in an oscillator.
\newblock \textit{Physical Review A}, 64(1):012310, 2001.

\bibitem{Hallgren2007}
S.~Hallgren.
\newblock Polynomial-time quantum algorithms for Pell's equation and the
  principal ideal problem.
\newblock \textit{J. ACM}, 54(1):4, 2007.

\bibitem{Harrington2001}
J.~Harrington and J.~Preskill.
\newblock Achievable rates for the Gaussian quantum channel.
\newblock \textit{Physical Review A}, 64:062301, 2001.

\bibitem{Kannan1983}
R.~Kannan.
\newblock Improved algorithms for integer programming and related lattice
  problems.
\newblock In \textit{Proc. STOC}, pages 193--206, 1983.

\bibitem{Khot2005}
S.~Khot.
\newblock Hardness of approximating the shortest vector problem in lattices.
\newblock \textit{J. ACM}, 52(5):789--808, 2005.

\bibitem{LenstraLenstraLovász1982}
A.~K.~Lenstra, H.~W.~Lenstra, and L.~Lovász.
\newblock Factoring polynomials with rational coefficients.
\newblock \textit{Mathematische Annalen}, 261:515--534, 1982.

\bibitem{LyubashevskyMicciancio2006}
V.~Lyubashevsky and D.~Micciancio.
\newblock Generalized compact knapsacks are collision resistant.
\newblock In \textit{ICALP}, Springer, 2006.

\bibitem{Lyubashevsky2022}
V.~Lyubashevsky et al.
\newblock Lattice-based group signatures and zero-knowledge proofs of
  automorphism stability.
\newblock In \textit{CCS}, 2022.

\bibitem{Menicucci2014}
N.~C.~Menicucci.
\newblock Fault-tolerant measurement-based quantum computing with
  continuous-variable cluster states.
\newblock \textit{Physical Review Letters}, 112(12):120504, 2014.

\bibitem{Micciancio2009}
D.~Micciancio and O.~Regev.
\newblock Lattice-based Cryptography.
\newblock In D.~Bernstein, J.~Buchmann, E.~Dahmen (eds.), \textit{Post-Quantum
  Cryptography}, Springer, 2009.

\bibitem{NielsenChuang2000}
M.~A.~Nielsen and I.~L.~Chuang.
\newblock \textit{Quantum Computation and Quantum Information}.
\newblock Cambridge University Press, 2000.

\bibitem{NIST2024}
{NIST}.
\newblock Post-Quantum Cryptography: Digital Signature Standard (FIPS 204, 205, 206).
\newblock \textit{Federal Information Processing Standards}, 2024.

\bibitem{Novaes2023}
N.~Novaes et al.
\newblock Machine-learning-guided lattice reduction.
\newblock \textit{IACR ePrint}, 2023.

\bibitem{PeikertRosen2007}
C.~Peikert and A.~Rosen.
\newblock Lattices that admit logarithmic worst-case to average-case connection
  factors.
\newblock In \textit{Proc. STOC}, 2007.

\bibitem{Regev2005}
O.~Regev.
\newblock On lattices, learning with errors, random linear codes, and
  cryptography.
\newblock In \textit{Proc. STOC}, pages 84--93, 2005.

\bibitem{Royer2022}
B.~Royer, S.~Singh, and S.~M.~Girvin.
\newblock Encoding qubits in multimode grid states.
\newblock \textit{PRX Quantum}, 3:010335, 2022.

\bibitem{Schnorr1994}
C.~P.~Schnorr.
\newblock Block Reduced Lattice Bases and Successive Minima.
\newblock \textit{Combinatorics, Probability and Computing}, 3:507--522, 1994.

\bibitem{Shor1994}
P.~W.~Shor.
\newblock Algorithms for quantum computation: Discrete logarithms and
  factoring.
\newblock In \textit{Proc. FOCS}, pages 124--134, 1994.

\bibitem{Sipser1996}
M.~Sipser and D.~Spielman.
\newblock Expander codes.
\newblock \textit{IEEE Trans. Inf. Theory}, 42(6):1710--1722, 1996.

\bibitem{Xu2019}
K.~Xu et al.
\newblock How powerful are graph neural networks?
\newblock In \textit{ICLR}, 2019.

\end{thebibliography}

\end{document}

